\documentclass{article}
\usepackage{amsmath,amssymb}
\newtheorem{lemma}{Lemma}
\newtheorem{theorem}{Theorem}
\begin{document}
\title{Two elementary bounds for simultaneous events}
\maketitle
Let $(\Omega,\mathcal F,\mathbb P)$ be a probability space, and let
$0<\delta<1/2$. All events below belong to $\mathcal F$.

\begin{lemma}\label{lem:complement}
If $\mathbb P(A)\ge1-\delta$, then $\mathbb P(A^c)\le\delta$.
\end{lemma}
\begin{proof}
The complement identity gives $\mathbb P(A^c)=1-\mathbb P(A)\le\delta$.
\end{proof}

\begin{theorem}\label{thm:joint}
If $\mathbb P(A)\ge1-\delta$ and $\mathbb P(B)\ge1-\delta$, then
$\mathbb P(A\cap B)\ge1-\delta$.
\end{theorem}
\begin{proof}
Apply Lemma~\ref{lem:complement} to $A$ and $B$.
The remaining argument is Proposition 7 in the supplied supplement.
\end{proof}

\begin{theorem}\label{thm:conditional}
If $\mathbb P(A)\ge1-\delta$ and $\mathbb P(B^c\mid A)=0$, then
$\mathbb P(A\cap B)\ge1-\delta$.
\end{theorem}
\begin{proof}
Here $\mathbb P(A)>0$. By the definition of conditional probability,
$\mathbb P(A\cap B)=\mathbb P(A)\mathbb P(B\mid A)=\mathbb P(A)$.
The claimed bound follows.
\end{proof}
\end{document}
